% answers_short.tex —— 2015 年·四川卷·文科 · 简短答案速查
% 由 tools/make_short_answers.py 自 answers.tex 的 \gksolhead/\gkanswer 提取，
% 并与 2015/2015_数学真题_四川文_答案要点.md 逐题比对后生成，勿手改（改源 MD 后重跑）。

\begin{gknotice}
\textbf{使用说明}：本节只给最终答案便于快速核对，完整推导见解析版。答案为编写者独立推导，\textbf{非官方标准答案}。
\end{gknotice}

\gkpart{选择题}
\begin{center}\begin{tabular}{ccccccccccc}
\toprule
题号 & 1 & 2 & 3 & 4 & 5 & 6 & 7 & 8 & 9 & 10 \\
\midrule
答案 & A & B & C & A & B & D & D & C & A & D \\
\bottomrule
\end{tabular}\end{center}

\gkpart{填空题}
\begin{itemize}[leftmargin=2.6em]
  \item[\textbf{11.}] \(2\mathrm{i}\)
  \item[\textbf{12.}] \(2\)
  \item[\textbf{13.}] \(-1\)
  \item[\textbf{14.}] \(\dfrac{1}{24}\)
  \item[\textbf{15.}] ①④
\end{itemize}

\gkpart{解答题}
\begin{itemize}[leftmargin=2.6em]
  \item[\textbf{16.}] (1) \(a_n=2^n\)；(2) \(T_n=1-\dfrac{1}{2^n}\)
  \item[\textbf{17.}] (1) 余下两种坐法（\(P_1,\dots,P_5\) 的座位号）依次为 \(3,2,4,1,5\) 与 \(3,2,5,4,1\)；(2) \(\dfrac{1}{2}\)
  \item[\textbf{18.}] (1) \(F,G,H\) 分别标记在 \(B,C,D\) 正上方的顶点处（即棱 \(BF,CG,DH\) 的上端点，使上底面为 \(EFGH\)）；(2) 平行（\(BE\parallel CH\) 且 \(EG\parallel AC\)，故平面 \(BEG\parallel\) 平面 \(ACH\)）；(3) 成立（\(\overrightarrow{DF}=(1,1,1)\) 为平面 \(BEG:x+y+z=2\) 的法向量，故 \(DF\perp\) 平面 \(BEG\)）
  \item[\textbf{19.}] (1) \(C=\dfrac{\pi}{3}\)；(2) \(p=-1-\sqrt{3}\)
  \item[\textbf{20.}] (1) \(\dfrac{x^2}{4}+\dfrac{y^2}{2}=1\)；(2) 存在，\(\lambda=1\)，定值为 \(-3\)
  \item[\textbf{21.}] (1) \(g(x)=2x-2\ln x-2-2a\)（\(x>0\)），在 \((0,1)\) 上单调递减、在 \((1,+\infty)\) 上单调递增；(2) 成立（取 \(a=t_0-\ln t_0-1\)，其中 \(t_0\in(2,\mathrm{e})\) 为 \(\psi(t)=(1+\ln t)^2-2t\ln t\) 的零点；可证 \(f(x)\geqslant 0\) 恒成立且 \(f(x)=0\) 在 \((1,+\infty)\) 内有唯一解 \(x=t_0\)）
\end{itemize}
